\documentclass[10pt,a4paper]{article}

\usepackage[top=1.7cm,bottom=1.8cm,left=1.8cm,right=1.8cm]{geometry}
\usepackage[T1]{fontenc}
\usepackage[utf8]{inputenc}
\usepackage{graphicx}
\usepackage{float}
\usepackage{xcolor}
\usepackage{caption}
\usepackage{booktabs}
\usepackage{array}
\usepackage{hyperref}
\usepackage{titlesec}

\definecolor{ink}{HTML}{252525}
\definecolor{muted}{HTML}{66645F}
\definecolor{accent}{HTML}{2878D0}
\definecolor{soft}{HTML}{F3F5F7}

\hypersetup{colorlinks=true,urlcolor=accent,linkcolor=accent}
\setlength{\parindent}{0pt}
\setlength{\parskip}{0.55em}
\setlength{\emergencystretch}{1em}
\captionsetup{font=footnotesize,labelfont=bf,skip=5pt}
\titleformat{\section}{\Large\bfseries\color{ink}}{\thesection}{0.7em}{}
\titleformat{\subsection}{\large\bfseries\color{ink}}{\thesubsection}{0.6em}{}

\newcommand{\paperquote}[1]{\textcolor{muted}{\textit{The atlas reports ``#1.''}}}
\newcommand{\evidence}[1]{\textcolor{muted}{\textit{Evidence strength: #1.}}\par}
\newcommand{\reportfigure}[4][0.94\textwidth]{%
  \begin{figure}[H]
    \centering
    \includegraphics[width=#1]{summary_downloads/#2}
    \caption{#3}
    \label{#4}
  \end{figure}
}

\begin{document}

\begin{center}
  {\LARGE\bfseries Higher-order association rules in gastrointestinal aGVHD}\\[0.35em]
  {\large Proof of added information, atlas alignment, and new biological hypotheses}
\end{center}
\vspace{0.25cm}
\hrule
\vspace{0.45cm}

\colorbox{soft}{\parbox{0.96\textwidth}{
\textbf{How to read this report.} A higher-order rule contains more than two cell types.
A rule is \emph{new} in one FOV when it passes rule-level FDR $\leq0.05$ but none of its
constituent pair rules was mined in that FOV. New attractions require Lift $\geq1.3$;
new avoidances require inverse Lift, $1/\mathrm{Lift}$, $\geq1.3$. Stage analyses require
at least 20 cells of every rule cell type and at least five rule-bearing FOVs in one stage.
For a new rule, the pair comparison is a lower bound based on the pair-mining threshold,
not a measured pair-to-complex ratio. Individual FOVs have rule-level statistical support;
stage prevalence and patient-median Lift are descriptive and are not patient-level
significance tests. Spatial association does not establish causality.
}}

% =============================================================================
\section{Proof that higher-order mining adds information}
% =============================================================================

The first test is global: after restricting to rules with FDR $\leq0.05$, does the method
mostly repeat shorter rules, or does it isolate a visible, structured remainder? Most long
rules are correctly labeled redundant, but the nonredundant fraction is substantial and
depends on rule shape. For attractions, 5--21\% are clearly stronger than a mined parent and
6--12\% have no mined parent. For avoidances, clearly stronger rules reach 27\% for
one-to-many rules. This separation is the method proof: the algorithm does not claim that
every longer rule is useful; it distinguishes inherited pair effects from joint contexts.

\evidence{Strong as an internal method check; these are rule counts, not biological effect sizes}

\reportfigure{complex_attracts_verdicts.png}{Verdicts among statistically supported
attractive rules. Percentages are within rule shape; totals at right show that the result is
not driven by a small rule family.}{fig:attract-verdicts}

\reportfigure{complex_avoids_verdicts.png}{The same classification for statistically
supported avoidance rules. The result confirms that both spatial directions contain
nonredundant higher-order information, while most longer rules are still rejected as
shorter-rule restatements.}{fig:avoid-verdicts}

\subsection{Two direct ways a complex rule can beat its simpler rule}

The strongest direct example is Endothelial, Epithelial $\rightarrow$ Goblet, SMV in a
Colon FOV. Its Lift is 23.3 versus 4.45 for the strongest simpler rule, a measured
$5.2\times$ gain. The second example is different: adding Epithelial to Endothelial,
Macrophage $\rightarrow$ CD8T raises Lift only from 1.34 to 1.45, but changes the FDR from
0.743 to 0.00189. Thus higher-order mining can add either a much larger effect or a joint
context in which an otherwise unsupported relationship becomes statistically detectable.

\evidence{Strong at the displayed FOV level; examples are selected automatically by the notebook}

\reportfigure{complex_attracts_stronger_fov.png}{A measured stronger-effect example. The
complete FOV, higher-order rule and strongest mined simpler rule use one color map and print
Lift, support, conviction and FDR where available.}{fig:direct-stronger}

\reportfigure{complex_attracts_significance_fov.png}{A measured significance-gain example.
The additional Epithelial context converts the strongest parent from unsupported to
rule-level FDR $<0.05$ while preserving a modest Lift gain.}{fig:direct-significance}

\subsection{New rules recur when no pair rule is available}

New rules provide the strictest proof of added representational capacity: no constituent
pair was mined in the same FOV. The strongest recurring attractive families include
epithelial--plasma, stromal--vascular and neural--epithelial contexts. Avoidances also recur,
although their lower-bound gains are generally smaller.

\begin{table}[H]
\centering
\footnotesize
\begin{tabular}{>{\raggedright\arraybackslash}p{7.0cm}rrr}
\toprule
New higher-order attraction & FOVs & median Lift & gain floor \\
\midrule
Epithelial, Macrophage $\rightarrow$ Plasma & 24 & 1.45 & $\geq1.21\times$ \\
Epithelial, SMV $\rightarrow$ Plasma & 22 & 1.43 & $\geq1.19\times$ \\
Epithelial, Fibroblast $\rightarrow$ Plasma & 18 & 1.43 & $\geq1.19\times$ \\
Epithelial, SMV $\rightarrow$ Endothelial & 13 & 2.76 & $\geq2.30\times$ \\
Epithelial, Neuron $\rightarrow$ Paneth & 12 & 2.24 & $\geq1.86\times$ \\
\bottomrule
\end{tabular}
\caption{Representative recurrent new attractions, ordered primarily by rule-bearing FOVs.
The gain floor compares the complex Lift with the largest value allowed for an unmined pair.}
\label{tab:new-attracts}
\end{table}

\begin{table}[H]
\centering
\footnotesize
\begin{tabular}{>{\raggedright\arraybackslash}p{7.0cm}rrr}
\toprule
New higher-order avoidance & FOVs & median inverse Lift & gain floor \\
\midrule
Muscle $\rightarrow$ Fibroblast, Macrophage & 17 & 1.40 & $\geq1.12\times$ \\
Fibroblast, Neuron $\rightarrow$ Goblet & 13 & 1.45 & $\geq1.16\times$ \\
Endothelial, Fibroblast $\rightarrow$ Goblet & 12 & 1.66 & $\geq1.33\times$ \\
Muscle, SMV $\rightarrow$ Fibroblast & 12 & 1.43 & $\geq1.15\times$ \\
Muscle, SMV $\rightarrow$ Macrophage & 11 & 1.44 & $\geq1.15\times$ \\
\bottomrule
\end{tabular}
\caption{Representative recurrent new avoidances. Inverse Lift puts attraction and
avoidance strength on the same ``larger is stronger'' scale.}
\label{tab:new-avoids}
\end{table}

\reportfigure{complex_new_gain_ranking.png}{The largest lower-bound gains among recurring
new attractions. Every rule passes FDR 0.05, Lift 1.3 and the five-FOV recurrence filter.}
{fig:new-attract-ranking}

\reportfigure{complex_new_avoid_gain_ranking.png}{The corresponding ranking for new
avoidances, using inverse Lift. The recurrent set is real but less strongly separated from
the pair-mining boundary than the leading attractions.}{fig:new-avoid-ranking}

\subsection{The stage result is not explained only by cell availability}

A non-new Colon rule provides a controlled comparison because its strongest simpler rule
is measurable. Goblet, Macrophage $\rightarrow$ Epithelial, Plasma falls from 48\% of
Control FOVs to 2\% of Mild and 9\% of Severe FOVs when all FOVs form the denominator.
After requiring 20 cells of every rule type, the complex prevalences are 48\%, 8\% and 21\%,
while the strongest simpler rule remains much more common at 92\%, 58\% and 79\%.
Eligibility changes the denominator but does not erase the loss of full joint organization.

\evidence{Descriptive availability control; not a patient-level stage test}

Figure~\ref{fig:availability-states} displays the missing control directly. Every stage FOV
is partitioned into too few cells to test, eligible but empty, covered by a simpler rule, or
carrying the informative complex rule. Thus the reader can see cell availability and spatial
organization in the same denominator rather than infer one from two separate percentages.

\reportfigure{complex_nonnew_goblet_macrophage_epithelial_plasma_availability.png}
{Every Colon FOV is counted exactly once. Grey is insufficient cell availability; the other
segments are eligible FOVs split by what spatial rule was found. The informative-rule segment
shrinks in disease even after the unavailable FOVs are made explicit.}{fig:availability-states}

\begin{figure}[H]
  \centering
  \includegraphics[width=0.92\textwidth]{summary_downloads/complex_nonnew_all_goblet_macrophage_epithelial_plasma_strength.png}
  \vspace{0.35em}

  \includegraphics[width=0.92\textwidth]{summary_downloads/complex_nonnew_eligible_goblet_macrophage_epithelial_plasma_strength.png}
  \caption{The same Colon rule before (upper) and after (lower) cell-eligibility filtering.
  Both prevalence lines compare the complex and strongest simpler rule; the right panels
  compare their patient-level Lift in rule-bearing FOVs.}
  \label{fig:availability-check}
\end{figure}

The aggregate comparison establishes the stage pattern; Figure~\ref{fig:availability-fovs}
then makes the selected tissue examples inspectable rather than asking the reader to trust
the summary alone.

\reportfigure{complex_nonnew_eligible_goblet_macrophage_epithelial_plasma_fovs.png}
{Eligible Control, Mild and Severe examples for the availability check. Each column shows
the full FOV, the complex rule and its strongest simpler rule.}{fig:availability-fovs}

\subsection{Complex and simpler rules can move in opposite directions}

The method should also expose cases in which the longer organization changes differently
from its parts. Figure~\ref{fig:opposite-movement} provides a filtered example:
Epithelial, Neuron $\rightarrow$ Goblet, Macrophage becomes more prevalent from Control to
Severe while the constituent simpler-rule coverage decreases slightly. In rule-bearing FOVs,
the complex Lift decreases from 2.86 to 2.43 while the simpler-rule Lift increases from 1.91
to 2.21. The two panels therefore show genuine opposite movement in both summaries, not a
duplicated Lift axis.

\evidence{Descriptive reversal passing the five-FOV recurrence filter; patient counts are printed in the figure}

\reportfigure{complex_opposite_epithelial_neuron_goblet_macrophage_strength.png}
{An explicit simple-versus-complex reversal. The left panel compares prevalence across the
same eligible stage FOVs; the right compares patient-median Lift only where each rule is
observed.}{fig:opposite-movement}

The same screen found no new-only rule with opposite complex-versus-simpler prevalence
movement between any two stages after the five-FOV filter. This is expected to be harder:
for a new occurrence, no constituent pair is mined in that same FOV, so the only possible
comparison is with pair-rule prevalence elsewhere in the stage. No weaker example is shown.

% =============================================================================
\clearpage
\section{Higher-order results aligned with the spatial-atlas paper}
% =============================================================================

\subsection{Loss of organized plasma-cell contexts}

The \href{https://doi.org/10.1126/scitranslmed.adu6032}{spatial atlas} reports that
``CD4 T cells, plasma cells, and macrophages were spatially organized from the crypts to the
tip of the villus'' and that homeostatic immune organization is disrupted with a prominent
reduction in IgA-secreting plasma cells. The complex-rule screen independently recovers a
family of Colon plasma contexts. Figure~\ref{fig:new-colon-overview} tests whether this is
one selected rule or a coordinated family; Figure~\ref{fig:plasma-availability} then tests
whether the leading rule disappears merely because its cell types become unavailable.

\evidence{Strong descriptive family-level alignment; stage contrasts are not permutation-tested here}

\reportfigure{complex_new_colon_attracts_overview.png}{All recurring new Colon attractions
that pass the current filters. Five of six rows end in Plasma and show a coordinated loss
from Control. Exact eligible denominators are printed in every cell.}{fig:new-colon-overview}

\reportfigure{complex_epithelial_macrophage_plasma_availability.png}
{Cell availability and spatial organization for the leading Colon plasma rule. Each stage
bar uses all FOVs: grey marks too few cells, while the plum segment marks eligible FOVs with
the new complex rule. The loss of the plum segment is visible separately from changing cell
availability.}{fig:plasma-availability}

Together, the figures show five related new rules ending in Plasma, with 20--24\% prevalence
in eligible Control FOVs but at most 8\% in Mild and 6\% in Severe. Epithelial, Macrophage
$\rightarrow$ Plasma is also the most recurrent new rule overall (24 FOVs; median Lift 1.45).
It refines rather than repeats the atlas result because no constituent pair was mined in
those rule-bearing FOVs. Figure~\ref{fig:epi-macro-plasma} shows the automatically selected
Colon examples behind that summary.

\reportfigure{complex_epithelial_macrophage_plasma_colon_fovs.png}{Stage-specific Colon FOV
examples for Epithelial, Macrophage $\rightarrow$ Plasma. Because the rule is new, each
available column shows the full FOV and the highlighted higher-order rule without inventing
a simpler comparator.}
{fig:epi-macro-plasma}

\subsection{Stromal--vascular remodeling is specifically higher-order}

The atlas describes increased fibrosis and broad stromal disruption. In eligible Colon
FOVs, the screen identifies Endothelial, Epithelial $\rightarrow$ Goblet, SMV as the clearest
Severe-enriched non-new attraction. The next three figures expose the complete argument:
its position in the filtered screen, the availability-versus-organization partition, and the
direct complex-versus-simpler prevalence and Lift comparison.

\evidence{Strong descriptive separation from the simpler rule; Severe includes five complex-rule patients}

\reportfigure{complex_eligible_colon_attracts_overview.png}{Top eligible non-new Colon
attractions. The stromal--vascular rule is the clearest Severe-enriched pattern; the same
overview also shows the Control-enriched plasma contexts.}{fig:eligible-colon-overview}

\reportfigure{complex_endothelial_epithelial_goblet_smv_availability.png}
{All Colon FOVs partitioned by whether the four required cell types are sufficiently
available and, if so, whether the FOV contains neither rule, only simpler coverage, or the
informative complex rule. The Severe enrichment is visible in the final segment and cannot
be read as cell availability alone.}{fig:endo-epi-availability}

\reportfigure{complex_eligible_endothelial_epithelial_goblet_smv_strength.png}{Prevalence and
Lift for Endothelial, Epithelial $\rightarrow$ Goblet, SMV and its strongest simpler rule.
The complex rule reaches 6/10 eligible Severe FOVs; patient counts are printed below the
Lift panel.}{fig:endo-epi-strength}

Figures~\ref{fig:endo-epi-availability}--\ref{fig:endo-epi-strength} show the numerical
separation: complex prevalence rises from 13\% and 11\% to 60\%, while simpler-rule coverage
changes in the opposite direction from 93\% to 90\%. The complex Lift also exceeds the
simpler Lift in the Control and Severe medians. The selected maps below show what the added
four-cell context looks like in tissue.

\reportfigure{complex_eligible_endothelial_epithelial_goblet_smv_fovs.png}{Automatically
selected examples show how the four-cell organization occupies tissue beyond the nearly
ubiquitous Endothelial $\rightarrow$ SMV parent.}{fig:endo-epi-fovs}

A separate parentless rule, Epithelial, SMV $\rightarrow$ Endothelial, supports the same
remodeling theme without depending on the non-new Colon result. Figure~\ref{fig:epi-smv-occurrences}
is a reproducibility check, not a disease-trend test. It shows every occurrence so that the
reader can see whether the result is a single-field accident: 13 FOVs from 13 patients carry
the rule, with median Lift 2.76 and a gain floor of $2.30\times$.

\reportfigure[0.82\textwidth]{complex_epithelial_smv_endothelial_occurrences.png}{Every FOV
in which Epithelial, SMV $\rightarrow$ Endothelial is classified as new, grouped by organ
and Pathological-score stage. This panel tests recurrence across independent patients; it
does not estimate a disease-stage trend. Diamonds are medians and the dashed line is the
predeclared Lift threshold.}
{fig:epi-smv-occurrences}

The recurrence plot establishes that this is not a single chosen field; the next panel
makes one automatically selected occurrence spatially inspectable.

\reportfigure[0.72\textwidth]{complex_epithelial_smv_endothelial_fov.png}{A representative
new Epithelial, SMV $\rightarrow$ Endothelial occurrence. No constituent pair was mined in
this FOV.}{fig:epi-smv-endothelial}

\subsection{A higher-order extension of Duodenal CD8 heterogeneity}

The atlas reports ``marked intrapatient heterogeneity in the prevalence of CD8 T cells and
neutrophils in the vicinity of individual crypts and villi.'' The pair-rule analysis found
increasing Duodenal CD8T--Goblet avoidance. Complex mining adds a macrophage-conditioned
version: Goblet $\rightarrow$ CD8T, Macrophage. The overview first places it in the complete
eligible avoidance screen; the availability and strength figures then display the exact
denominators, pair coverage and Lift trajectory needed to judge the claim.

\evidence{Moderate paper alignment; good FOV denominators but only two Severe patients}

\reportfigure{complex_eligible_duodenum_avoids_overview.png}{The eligible Duodenal avoidance
screen. Multiple CD8T-containing and muscle-containing higher-order rules become more
prevalent in Severe disease. Rows with one or two eligible Severe FOVs are visible and are
not used as primary claims.}{fig:eligible-duodenum-avoid-overview}

\reportfigure{complex_goblet_cd8t_macrophage_availability.png}
{Every Duodenal FOV is assigned to one availability/organization state for Goblet
$\rightarrow$ CD8T, Macrophage. This shows the changing supply of measurable FOVs separately
from the fraction carrying the informative higher-order avoidance.}{fig:goblet-cd8-availability}

\reportfigure{complex_eligible_goblet_avoids_cd8t_macrophage_strength.png}{Prevalence and
Lift of Goblet $\rightarrow$ CD8T, Macrophage against its strongest simpler rule. For
avoidance, lower Lift means stronger separation.}{fig:goblet-cd8-macro-strength}

The figures show complex prevalence rising from 9/28 eligible Control FOVs to 26/50 Mild
and 7/8 Severe FOVs, while simpler-rule coverage is already 89--100\%. Median complex Lift
falls from about 0.42 to 0.29, so the avoidance strengthens numerically. Only two Severe
patients contribute to that Lift estimate, so it remains descriptive. The following maps
show the actual joint spatial pattern used for the summary.

\reportfigure{complex_eligible_goblet_avoids_cd8t_macrophage_fovs.png}{Control, Mild and
Severe examples of the macrophage-conditioned CD8T avoidance. These maps make the joint
context inspectable but do not identify the atlas's patient-level CD8 subgroup.}
{fig:goblet-cd8-macro-fovs}

% =============================================================================
\clearpage
\section{New biologically interpretable findings}
% =============================================================================

The complete Duodenal new-attraction screen organizes the selected hypotheses into three
recurring themes: an epithelial--neural Paneth context, several epithelial--plasma contexts,
and smooth-muscle/mesenchymal contexts ending in Endothelial. The overview is shown before
the individual rules so that the examples can be judged against the full filtered result
rather than as isolated selections.

\reportfigure{complex_new_duodenum_attracts_overview.png}{All recurring new Duodenal
attractions that pass rule-level FDR, Lift, cell-eligibility and recurrence filters. Exact
eligible denominators expose which apparent stage differences can and cannot be trusted.}
{fig:new-duodenum-attract-overview}

\subsection{A surviving epithelial--neural Paneth niche}

Epithelial, Neuron $\rightarrow$ Paneth is new in 12 FOVs, with median Lift 2.24 and a gain
floor of $1.86\times$. This extends the atlas observation of Paneth-cell loss: among FOVs
that retain enough Paneth cells to measure, Paneth proximity is associated with a joint
epithelial--neural context that no simpler mined rule captures. Neuronal signaling can
regulate crypt growth, and receptors for that signaling occur in the crypt stem-cell
compartment (\href{https://pubmed.ncbi.nlm.nih.gov/28054012/}{Muise et al., 2017}), making
the niche plausible without establishing direct Neuron--Paneth signaling.

\evidence{Moderate for the recurring higher-order niche; exploratory mechanism}

\reportfigure[0.80\textwidth]{complex_epithelial_neuron_paneth_strength.png}{The rule occurs
in 3/11 eligible Control and 8/58 Mild FOVs, representing two and seven patients. No Severe
rule-bearing FOV passes the filters, so the line is not a validated severity trend.}
{fig:neuron-paneth-strength}

The prevalence panel establishes recurrence but not morphology. The next figure shows the
automatically selected Control and Mild fields, with all non-rule cell types greyed in the
highlighted row.

\reportfigure{complex_epithelial_neuron_paneth_fovs.png}{Representative Control and Mild
FOVs for the new Paneth rule. The empty Severe column makes the absence of a qualifying
example explicit.}{fig:neuron-paneth-fovs}

\subsection{A recurrent neuro--stromal context in Colon}

Fibroblast, Muscle $\rightarrow$ Neuron is a new rule in eight Colon FOVs. It reaches 5/15
eligible Mild FOVs and is also observed in Control and Severe; median Lift is 2.23, 2.67 and
1.50 respectively. The result is biologically coherent with the close organization of
enteric nerves and the muscular--stromal compartment, but the patient counts are one, four
and two. The robust conclusion is recurrence of a three-cell context, not a stage trend.

\evidence{Moderate recurrence; exploratory stage interpretation}

\reportfigure[0.82\textwidth]{complex_fibroblast_muscle_neuron_strength.png}{Prevalence and
Lift for the new Colon neuro--stromal rule. Exact FOV and patient counts prevent the Mild
peak from being overinterpreted.}{fig:fibro-muscle-neuron-strength}

To verify that the numerical recurrence corresponds to coherent tissue structure, the next
panel displays the selected rule-bearing FOV from every stage in which one is available.

\reportfigure{complex_fibroblast_muscle_neuron_fovs.png}{Representative rule-bearing FOVs
from all three stages. The highlighted maps reveal a spatially coherent muscular boundary
in the selected Control and Mild examples.}{fig:fibro-muscle-neuron-fovs}

\subsection{Disease-associated niches among surviving plasma and endothelial cells}

Epithelial, SMV $\rightarrow$ Plasma\_CD38 is absent from 26 eligible Control FOVs but
present in 7/70 Mild and 3/22 Severe FOVs, with Lift above 1.3. This does not contradict the
atlas's overall plasma-cell loss: it suggests that when CD38-positive plasma cells remain,
some occupy a disease-associated epithelial--SMV context. Five Mild and two Severe patients
carry the rule, so it is a follow-up hypothesis rather than a validated disease trajectory.

\evidence{Moderate recurrence with sound eligible denominators; limited patient replication}

\reportfigure[0.82\textwidth]{complex_epithelial_smv_plasma_cd38_strength.png}{The new
Plasma\_CD38 context is not observed in eligible Control FOVs and has median Lift 1.50 in
Mild and 1.67 in Severe rule-bearing patients.}{fig:plasma-cd38-strength}

The next maps show the Mild and Severe occurrences used to ground that stage-level summary;
the missing Control example is part of the result, not a plotting omission.

\reportfigure{complex_epithelial_smv_plasma_cd38_fovs.png}{Representative Mild and Severe
occurrences. No Control FOV carries the rule despite 26 eligible FOVs.}
{fig:plasma-cd38-fovs}

Mesenchymal\_VIM, SMV $\rightarrow$ Endothelial is detected in five Mild Duodenal FOVs,
each from a different patient, with median Lift 2.31. The atlas's fibrosis and stromal
disruption make the relationship biologically plausible, but only one Control FOV is
eligible. It is therefore evidence for a recurrent Mild context, not evidence for a
Control-to-Mild increase.

\evidence{Exploratory; five-patient recurrence but an inadequate Control denominator}

\reportfigure[0.82\textwidth]{complex_mesenchymal_smv_endothelial_strength.png}{The Mild-only
pattern and its eligibility limitation. Control prevalence is 0/1, so the apparent rise
must not be interpreted as a stage comparison.}{fig:mesenchymal-endo-strength}

Because the Control denominator prevents a stage claim, the following panel is presented
only as spatial confirmation of the recurrent Mild configuration.

\reportfigure[0.80\textwidth]{complex_mesenchymal_smv_endothelial_fovs.png}{A representative
Mild FOV containing the new mesenchymal--smooth-muscle--endothelial context.}
{fig:mesenchymal-endo-fovs}

\subsection{Negative rules: higher-order spatial avoidance}

Avoidance is analyzed as a complete result family, using the same FDR, cell-eligibility,
strength and recurrence filters as attraction. No new Colon avoidance reaches five FOVs in
one eligible stage. The retained Duodenal screen is small and coherent: all five rules
contain Muscle or SMV, but no stage prevalence exceeds 11\%. This supports a recurring
muscle-centered family while arguing against a strong new-avoidance stage claim.

\reportfigure{complex_new_duodenum_avoids_overview.png}{All recurring new Duodenal
avoidances. The common muscle/SMV theme is biologically interesting, but the low prevalence
keeps the stage interpretation exploratory.}{fig:new-duodenum-avoid-overview}

The new avoidance Muscle $\rightarrow$ Fibroblast, Macrophage recurs in 17 FOVs, with
median inverse Lift 1.40. A representative Duodenal FOV has Lift 0.717, support 0.0585,
conviction 0.576 and FDR 0.00656. This joins the Fibroblast/Muscle/Neuron finding in pointing
to the muscular--stromal compartment, but in the opposite spatial direction. Muscularis
macrophages participate in gut neuromuscular homeostasis
(\href{https://pmc.ncbi.nlm.nih.gov/articles/PMC4149228/}{M\"uller et al., 2014}); the
present rule does not identify that macrophage subtype and should be followed up by cell
state or marker analysis.

\evidence{Moderate recurrence; biologically plausible but subtype-agnostic}

\reportfigure[0.82\textwidth]{complex_muscle_avoids_fibroblast_macrophage_strength.png}
{Eligible stage prevalence and Lift for the new Muscle $\rightarrow$ Fibroblast, Macrophage
avoidance. The exact FOV and patient counts show both its recurrence and why its stage trend
remains exploratory.}{fig:muscle-fibro-macro-strength}

The recurrence statistics motivate inspection of stage-specific fields. Since this is a
new rule, the spatial panel correctly contains no simpler-rule row.

\reportfigure{complex_muscle_avoids_fibroblast_macrophage_stage_fovs.png}{Automatically
selected stage examples of the new Muscle $\rightarrow$ Fibroblast, Macrophage avoidance.
Each available column shows the full field and the higher-order rule.}
{fig:muscle-fibro-macro-fov}

A complementary non-new avoidance, Muscle $\rightarrow$ Plasma, SMV, becomes substantially
more prevalent in disease while its simpler-rule coverage remains high. Its complex Lift
rises from 0.04 in Control to 0.17 in Mild and 0.11 in Severe: the avoidance is much more
frequent in disease but is not stronger than the single Control estimate. The next figure
shows both summaries and both rules, preventing prevalence from being confused with effect
strength. Only one Control patient carries the complex rule, so the Lift trajectory is
descriptive.

\evidence{Strong FOV-level prevalence contrast; Control strength is based on one patient}

\reportfigure{complex_eligible_muscle_avoids_plasma_smv_strength.png}{The non-new Duodenal
Muscle $\rightarrow$ Plasma, SMV avoidance rises from 1/17 eligible Control FOVs to 19/58
Mild and 9/15 Severe FOVs; its strongest simpler rule is shown in both panels.}
{fig:muscle-plasma-smv-strength}

A distinct non-new rule, Muscle $\rightarrow$ CD8T, Fibroblast, has adequate eligible
denominators in every stage. Complex-rule prevalence rises from 3/17 Control and 5/36 Mild
FOVs to 13/21 Severe FOVs, whereas the strongest simpler rule is more common throughout.
At the same time, complex Lift rises from 0.20 to 0.45, meaning the avoidance becomes weaker
among rule-bearing FOVs. This divergence is important: prevalence and effect strength are
different biological summaries and neither should substitute for the other.

\evidence{Strong descriptive example of prevalence--strength divergence; six Severe patients}

\reportfigure{complex_eligible_muscle_avoids_cd8t_fibroblast_strength.png}{The complex rule
becomes more prevalent but less strongly avoidant across stages. The side-by-side panels
prevent the prevalence increase from being misread as stronger spatial separation.}
{fig:muscle-cd8-fibro-strength}

The final maps show the stage examples behind both lines and retain the strongest simpler
rule so the reader can inspect what the extra Fibroblast context adds.

\reportfigure{complex_eligible_muscle_avoids_cd8t_fibroblast_fovs.png}{Representative
Control, Mild and Severe examples of the complex and strongest simpler rule.}
{fig:muscle-cd8-fibro-fovs}

\end{document}
